\documentclass[11pt,letterpaper]{article}
\usepackage[T1]{fontenc}
\usepackage{lmodern,amsmath,amssymb,amsthm,mathtools,booktabs,microtype,xcolor}
\usepackage[margin=.88in,headheight=15pt]{geometry}
\usepackage{fancyhdr,xurl,hyperref}
\definecolor{navy}{HTML}{19354A}
\definecolor{teal}{HTML}{126C70}
\hypersetup{colorlinks=true,linkcolor=navy,citecolor=teal,urlcolor=teal,pdftitle={Polynomial auxiliary scales and Kummer prime races beyond Dirichlet GRH}}
\pagestyle{fancy}\fancyhf{}\lhead{\small\color{navy} GRH method transfer}\rhead{\small\color{navy} 25 November 2026}\cfoot{\thepage}
\setlength{\parindent}{0pt}\setlength{\parskip}{5pt plus 1pt}\setlength{\emergencystretch}{2em}
\newtheorem{theorem}{Theorem}[section]\newtheorem{lemma}[theorem]{Lemma}\newtheorem{proposition}[theorem]{Proposition}\newtheorem{corollary}[theorem]{Corollary}
\theoremstyle{definition}\newtheorem{definition}[theorem]{Definition}
\theoremstyle{remark}\newtheorem{remark}[theorem]{Remark}
\newcommand{\Q}{\mathbb Q}\newcommand{\C}{\mathbb C}\newcommand{\R}{\mathbb R}\newcommand{\F}{\mathbb F}\newcommand{\ord}{\operatorname{ord}}\newcommand{\Li}{\operatorname{Li}}\newcommand{\Gal}{\operatorname{Gal}}\newcommand{\Res}{\operatorname{Res}}\newcommand{\ind}{\operatorname{Ind}}\newcommand{\one}{\mathbf 1}
\title{\color{navy}\bfseries Polynomial auxiliary scales and\\ Kummer prime races beyond Dirichlet GRH\\[4pt]\large A sharper analytic transfer and an Artin primitive-root target}
\author{Research note}\date{25 November 2026 (Pacific date)}
\begin{document}\maketitle
\begin{abstract}
We examine the compact edition of the Alp\"oge--Claude Dirichlet GRH proof in \texttt{lsernau/grh-lean}, at tree/commit identifier \texttt{72fead102f663d5c7a407964a13dc0cf9d977207}. Two rigorously proved extensions are developed. First, a gamma-moment argument sharpens the source's auxiliary logarithmic interval $[t,t^2]$ to $[t,Ct]$, even for exponentially growing coefficients, and gives a common convexified scale for finitely many errors. Second, we prove a one-sided, fixed-log-window, positive-block criterion for simultaneous Dedekind RH for prime-index Kummer fields $\Q(\zeta_\ell,2^{1/\ell})$. Assuming the exported Dirichlet GRH, it can be expressed entirely as a deficit of nonresidue primes. These criteria, for blocks covering all odd primes, imply the Artin--Hooley density for the primitive root $2$; we give the sieve argument using only prime-index RH, not RH for all squarefree indices. A precise sufficient replacement for the source covolume inequality is isolated. That replacement is not proved here, so neither nonabelian RH nor Artin's conjecture is claimed solved. The analytic improvements are candidates for further research, not verified priority claims.
\end{abstract}
\section{Scope and relationship to previous work}
The exported theorem \texttt{LFunction\_ne\_zero\_of\_half\_lt\_re} is about \emph{Dirichlet characters modulo positive integers}. The four other exported statements concern RH and the location of zeros of Dirichlet $L$-functions. In particular, the endpoint is not RH for arbitrary Dedekind zeta functions, Hecke characters over arbitrary number fields, or automorphic $L$-functions. The repository's own README reports kernel and Comparator verification with the three standard Lean axioms; it also states that no mathematician has read through the whole proof. We take the stated Dirichlet result as the premise of the project, and independently prove the classical implications below. We inspected source files through GitHub, but did not rebuild the Lean project or replay its kernel. No independent certification of the upstream proof is asserted.

\paragraph{What is new in this packet.} Rather than enumerate familiar conditional consequences, we target improvements to the machinery itself. The first is a tighter, tunable scale-selection theorem that does not use the target GRH. The second is a combined test whose inputs are prime-index Kummer nonresidue races over $\Q$ and whose conclusion would give a specific instance of the still unresolved primitive-root problem. A block can contain arbitrarily many fixed primes; a single inequality per block suffices. The prime-only index restriction matters in the qualitative asymptotic sieve: finite inclusion-exclusion needs only unconditional Chebotarev, whereas the growing-index tail needs effective estimates. This last observation is a variant of Hooley's method, not a claimed discovery of a new sieve.

The incoming directory of ProveIt was listed using GitHub. Seven smaller archives were decoded in the connector session and selected README, report or evidence text was inspected. See \texttt{provenance.md} for the actual scope. Earlier reports already cover general one-sided zero criteria, Sylow and relative-tower compression, periodic Frobenius obstructions, fourth-power sampling, and some Kummer targets. We deliberately focus on the improved $[t,Ct]$ window, convexified multiple tests, fixed-log-window block averaging, and the explicit prime-index sieve for $2$. We cannot certify that none of these variants is already in an uninspected report or in the mathematical literature. We do not count the known consequences of Dirichlet GRH as new results.

\section{The transferable source architecture}\label{sec:source}
In the compact edition, \path{src/New/Main.lean} gives the central comparison. Let $D$ denote the admissible data, $d$ the target row dimension, $m$ the number of theta lifts, and $H$ the main height. The short/enough forms supply a family of almost full dimension with an archimedean bound and an index bound; integral forms supply a lower bound on the same covolume. The main height cancels. In the actual notation the result is
\begin{equation}\label{eq:star}D.m\,D.remainder\le (D.logB)^2 D.budget,\qquad D.budget=f a(\log b)^3.\end{equation}
The input primes are weighted by $\Lambda(n)(1-\chi_d(n))$ in a congruence class, with $\chi_d$ quadratic. The prime error has a ramp kernel $1/n-1/x$ and a logarithmic main term and constant. In \path{src/New/Endgame/PrimeError.lean}, \texttt{isPolar\_primeError} identifies a polar Laplace transform. In \path{src/New/Endgame/OneSidedCriterion.lean}, the declaration \texttt{IsPolar.exists\_window\_sub\_mul\_le} selects a good second scale in $[t,t^2]$ for $E(u)-X E'(u)$, $0\le X\le e^t$. The same file's \texttt{entireL\_ne\_zero\_of\_primeError\_le} uses an auxiliary quadratic twist to avoid cancellation of residues and a Landau argument to exclude zeros. Finally \texttt{eventually\_primeError\_le\_rpow} uses the saving $m\asymp\sqrt{x}$ to obtain the one-sided error $x^{-1/2+8\varepsilon}$.

\begin{center}\small
\begin{tabular}{p{.43\linewidth}p{.49\linewidth}}\toprule
Source declaration or mechanism & Role in the present transfer\\\midrule
\texttt{star\_inequality} & Model for the new, still missing arithmetic inequality.\\
\texttt{IsPolar.eventually\_moment\_le} & Cauchy bounds for gamma moments; proved here with elementary complex analysis.\\
\texttt{IsPolar.exists\_window\_sub\_mul\_le} & Improved logarithmic interval and vector extension (Section~\ref{sec:window}).\\
\texttt{isPolar\_primeError} & Model for ramp-prime Mellin identity, including centering constant.\\
\texttt{entireL\_ne\_zero\_of\_primeError\_le} & One-sided Landau mechanism; in our block the signs of residues make auxiliary twist nonvanishing unnecessary.\\\bottomrule
\end{tabular}\end{center}
The hard arithmetic modules are not interchangeable black boxes. The source constructs a quadratic CM setting, a tower of unitary groups and theta lifts, many short forms, index control, and fitted twists with local Selmer congruences. Changing the Frobenius weight changes the divisibility and height calculations. The mathematical results below do not assume that those constructions already work in nonabelian Kummer fields.

\section{A polynomial-size second scale}\label{sec:window}
\begin{theorem}[Tunable compact window]\label{thm:window}
Let $E_1,\ldots,E_d:[0,\infty)\to\R$ be continuous with $|E_i(v)|\le B e^v$. Assume that their Laplace transforms $H_i(z)=\int_0^\infty E_i(v)e^{-zv}\,dv$ (initially $\Re z>1$) admit holomorphic continuation to an open neighbourhood of the closed disk $|z-3|\le3$. Fix $\kappa,\lambda\ge0$. There is a constant $C_0>1$ such that for all sufficiently large $t$ and every real $a=(a_1,\ldots,a_d)$ with $\|a\|_1\le e^{\kappa t}$ there is $u\in[t,C_0t]$ with
\[
\left|\sum_{i=1}^d a_i E_i(u)\right|\le e^{-\lambda t}.
\]
Furthermore there is a probability distribution supported on at most $d+1$ points of $[t,C_0t]$, independent of $a$, for which every such linear combination has a mean whose absolute value is at most $e^{-\lambda t}$ (not a bound on the mean absolute value).
\end{theorem}
\begin{proof}
By compactness there is a common $R>3$ on whose closed disk all continuations are holomorphic in a neighbourhood; set $M=\max_i\max_{|z-3|=R}|H_i(z)|$. Cauchy's formula and differentiation under the integral at $3$ yield
\[
\left|\int_0^\infty v^k e^{-3v}E_i(v)\,dv\right|\le M k! R^{-k}.
\]
Choose
\begin{align*}
c&>\max\left\{\frac{\kappa+\lambda+1}{\log(R/3)},\frac{\kappa+\lambda+5}{\log2}\right\},\\
C_0&>\max\{1,\ \kappa+\lambda+1+c\log3\},\qquad k=\lceil ct\rceil.
\end{align*}
Let $V$ have gamma density $3^{k+1}v^ke^{-3v}/k!$ on $v\ge0$. Thus $|\mathbb E E_i(V)|\le 3M(3/R)^k$. Its moment generating function is $\mathbb E e^{sV}=(3/(3-s))^{k+1}$ for $s<3$. On the lower tail $V<t$, Markov's inequality gives
\[
\mathbb E[e^V\one_{V<t}]\le e^t e^{3t}\mathbb E e^{-3V}= e^{4t}2^{-(k+1)}.
\]
On the upper tail $V>C_0t$, we use $e^V\le e^{-C_0t}e^{2V}$ and get
\[
\mathbb E[e^V\one_{V>C_0t}]\le e^{-C_0t}3^{k+1}.
\]
Each of the moment and two tail bounds is $o(e^{-(\kappa+\lambda)t})$ by the strict choices of $c,C_0$. In particular $\mathbb P(t\le V\le C_0t)=1-o(1)$. The conditional mean vector $b=(\mathbb E[E_i(V)\mid t\le V\le C_0t])_i$ satisfies $\|b\|_\infty=o(e^{-(\kappa+\lambda)t})$. Hence $|a\cdot b|\le e^{-\lambda t}$ for every admissible $a$, eventually, uniformly in $a$. A continuous real function on a compact interval attains its conditional mean (between minimum and maximum), proving the single-scale claim. Finally, the continuous vector image is compact, its convex hull is compact in $\R^d$, and the conditional mean belongs to that convex hull by integration or separation. Carath\'eodory's theorem represents it by a convex combination of at most $d+1$ vector values. This is the required distribution.
\end{proof}

\begin{corollary}[Source-compatible form]\label{cor:window}
If $E,E'$ satisfy the hypotheses, then for large $t$ and every $0\le X\le e^t$ there is $u\in[t,C_0t]$ with $|E(u)-XE'(u)|\le e^{-t}$. In terms of $x_1=e^u$, taking $t=A\log x$ gives $x^A\le x_1\le x^{AC_0}$, with error at most $x^{-A}$.
\end{corollary}
\begin{proof}
For large $t$, $1+X\le e^{2t}$. Apply Theorem~\ref{thm:window} with $\kappa=2$, $\lambda=1$. The source's polar transforms are holomorphic near the closed disk in the shifted variable $z=s-1$: the disk lies in $\Re s\ge1$, where the continued centered transform is holomorphic. The source growth bound supplies $B$. Exponentiate the interval.
\end{proof}
\paragraph{Why the tightening matters.} The source allows $\log x_1$ up to a quadratic in $\log x$ in its scale selection. A fixed polynomial in $x$ now suffices. This does not improve the final $1/2$ exponent by itself: that exponent is the saving from the number of theta lifts. It does reduce the range of auxiliary levels that a new geometric construction would have to support. The constant $C_0$ depends on the analytic family through $R$; no uniformity in unbounded conductors is claimed. The vector statement does not assert a single common point for all errors (for example, a continuous curve can encircle zero while avoiding it). It yields a common weighted average, useful only where the arithmetic inequality is stable under averaging.

\section{Prime-index Kummer fields and the exact detector}\label{sec:kummer}
Let $\ell$ be an odd prime and define
\[
F_\ell=\Q(\zeta_\ell),\qquad B_\ell=\Q(2^{1/\ell}),\qquad K_\ell=F_\ell B_\ell.
\]
The degree of $B_\ell$ is $\ell$ by Eisenstein at $2$, and the degrees $\ell$ and $\ell-1$ are coprime. Thus $[K_\ell:\Q]=\ell(\ell-1)$, and the Galois group acts as all affine maps $x\mapsto ax+b$ of $\F_\ell$. Let $V_\ell$ denote the permutation character of this action minus the trivial character. Its values are $\ell-1$ at the identity, $-1$ at nonidentity translations, and $0$ when $a\ne1$. It is irreducible: restricting to the translations gives the direct sum of their nontrivial linear characters, and the multiplicative subgroup permutes these transitively. More precisely $V_\ell=\ind_N^G\psi$ for any nontrivial character $\psi$ of the translations $N$, since the induced character agrees with the listed values. Artin induction and abelian class field theory give
\begin{equation}\label{eq:fact}L_\ell(s)=L(s,V_\ell)=L_{F_\ell}(s,\psi)=\frac{\zeta_{B_\ell}(s)}{\zeta(s)},\qquad \zeta_{K_\ell}(s)=\zeta_{F_\ell}(s)L_\ell(s)^{\ell-1}.
\end{equation}
Indeed the regular character of $G$ is the sum of the inflated regular character of $G/N$ and $(\ell-1)V_\ell$; the permutation character for $B_\ell$ is $1+V_\ell$. The finite-order, nontrivial Hecke $L$-function $L_\ell$ is entire and nonzero on $\Re s\ge1$. Its zeros in the critical strip are symmetric about the critical line (the Artin character is real; in any case the standard functional equation and conjugation suffice). These classical facts do not assume nonabelian GRH.

For a prime $p\nmid 2\ell$, define $I_\ell(p)=\one_{p\equiv1\ (\mathrm{mod}\ \ell)}$ and $g_\ell(p)=\one_{p\equiv1\ (\mathrm{mod}\ \ell),\,2^{(p-1)/\ell}\equiv1\ (\mathrm{mod}\ p)}$. Let $f_\ell=I_\ell-g_\ell$ be the nonresidue indicator. All three weights are set to zero at the two bad primes. The complete-splitting condition is exactly $g_\ell=1$; one can equivalently count roots of $X^\ell-2$ over $\F_p$. In fact
\begin{equation}\label{eq:indicator}g_\ell=\frac{I_\ell+V_\ell}{\ell},\qquad f_\ell=\frac{(\ell-1)I_\ell-V_\ell}{\ell},\end{equation}
where $V_\ell$ means its value on Frobenius. This follows from the three affine possibilities above. The root-count test $r_\ell(p)-1=V_\ell(\operatorname{Frob}_p)$ is particularly cheap to compute without finding a splitting field.

For a real prime weight $w$ that is one of these finite group indicators (or their linear combinations), put
\begin{align*}
P_w(s)&=\sum_p w(p)\log p\,p^{-s},& W_w(x)&=\sum_{p\le x}w(p)\log p\left(\frac1p-\frac1x\right),\\
\delta_w&=\Res_{s=1}P_w(s),& c_w&=\lim_{s\to1}\left(P_w(s)-\frac{\delta_w}{s-1}\right)-\delta_w,\\
R_w(v)&=W_w(e^v)-\delta_w v-c_w.&&
\end{align*}
The continuations and constants exist by the identities below. The weights use primes only, rather than prime powers; this makes the input congruences and numerical checks simple. The difference is analytic for $\Re s>1/2$. In the half-plane of absolute convergence, orthogonality of the Dirichlet characters and \eqref{eq:indicator} give
\begin{align}
P_{g_\ell}(s)&=-\frac1{\ell(\ell-1)}\frac{\zeta_{F_\ell}'(s)}{\zeta_{F_\ell}(s)}-\frac1\ell\frac{L_\ell'(s)}{L_\ell(s)}+A_\ell(s),\label{eq:p-g}\\
P_{f_\ell}(s)&=-\frac1\ell\frac{\zeta_{F_\ell}'(s)}{\zeta_{F_\ell}(s)}+\frac1\ell\frac{L_\ell'(s)}{L_\ell(s)}+A'_\ell(s),\label{eq:p-f}
\end{align}
where $A_\ell,A'_\ell$ are holomorphic on $\Re s>1/2$. To verify this, logarithmically differentiate the Artin Euler factors: for unramified $p$ the coefficient at $p^j$ is $\chi(\operatorname{Frob}_p^j)\log p$. The $j\ge2$ terms converge absolutely locally uniformly for $\Re s>1/2$ because the traces are bounded for fixed degree; each bad-prime local factor is nonzero on $\Re s>0$ since the eigenvalues on inertia invariants are roots of unity. Removing these terms yields the stated holomorphic corrections. The same argument for the abelian regular character identifies $P_{I_\ell}$ with $-(\ell-1)^{-1}\zeta'_{F_\ell}/\zeta_{F_\ell}$ modulo a holomorphic function. The residues at $1$ are $\delta_g=1/(\ell(\ell-1))$, $\delta_f=1/\ell$, and $\delta_I=1/(\ell-1)$.

\begin{lemma}[Ramp transform and averaging]\label{lem:transform}
For $\Re s>1$, the Laplace transform in the shifted variable is
\begin{equation}\label{eq:transform}
\mathcal R_w(s)=\int_0^\infty R_w(v)e^{-(s-1)v}\,dv=\frac{P_w(s)}{s(s-1)}-\frac{\delta_w}{(s-1)^2}-\frac{c_w}{s-1}.
\end{equation}
It extends meromorphically to $\Re s>1/2$, holomorphically near $\Re s\ge1$. For $h>0$, define $\overline R_w(v)=h^{-1}\int_0^hR_w(v+u)\,du$. The corresponding transform is $J_h(s)\mathcal R_w(s)$ plus an entire function, with
\[
J_h(s)=\frac{e^{h(s-1)}-1}{h(s-1)},\qquad J_h(1)=1.
\]
This multiplier has no zeros for $\Re s<1$ and is positive for real $s$.
\end{lemma}
\begin{proof}
The integral for a single $p$ is $p^{-s}(1/(s-1)-1/s)$, multiplied by $w(p)\log p$. Absolute convergence justifies summation. The subtraction cancels the double and simple terms at $s=1$ by the definition of $c_w$; equations \eqref{eq:p-g}--\eqref{eq:p-f} and classical nonvanishing for the factors on $\Re s\ge1$ give continuation. For the moving average, replace $v$ by $v+u$ in the inner integral. The omitted initial interval $[0,u]$ contributes an entire function after integration in $u$ over $[0,h]$. The multiplier is the average of $e^{u(s-1)}$. Its only zeros are $s=1+2\pi i k/h$, $k\ne0$. Positivity on the real axis follows from its integral formula.
\end{proof}
All $R_w$ and their moving averages are continuous and $O((1+v)^2)$ by the crude bound $\sum_{n\le e^v}(\log n)/n=O((1+v)^2)$. They therefore satisfy the hypotheses of Theorem~\ref{thm:window} without invoking target RH.

\section{One averaged inequality per block}\label{sec:block}
\begin{lemma}[One-sided pole detector]\label{lem:landau}
Let $T$ be continuous, $T(v)=O((1+v)^2)$ on $[0,\infty)$. Suppose its shifted Laplace transform is meromorphic for $\Re s>\beta$, where $1/2<\beta<1$, regular for $\Re s\ge1$, and all its possible real poles in $(\beta,1)$ are simple with strictly positive residues. If $T(v)\le C e^{(\beta-1)v}$ eventually, then its transform is holomorphic for $\Re s>\beta$.
\end{lemma}
\begin{proof}
On a tail $v\ge v_0$ the function $U(v)=C e^{(\beta-1)v}-T(v)$ is nonnegative. Its shifted Laplace transform is $C/(s-\beta)-\mathcal T(s)$ plus an entire function. The real abscissa of convergence $\sigma_c$ is at most $1$. If $\sigma_c>\beta$, Landau's theorem for a Laplace transform of a nonnegative function gives a singularity at $\sigma_c$. For completeness: if the transform were analytic at $\sigma_c$, choose real $s_0>\sigma_c$ sufficiently close that its Taylor disk includes a real $s_0-r<\sigma_c$. The Taylor coefficients of $\mathcal U(s_0-z)$ are the nonnegative integrals of $U(v)e^{-(s_0-1)v}v^k/k!$. Monotone convergence at $z=r$ then shows convergence at $s_0-r$, a contradiction. The proposed singularity is impossible at $1$ by regularity, and elsewhere it would be a simple pole of negative residue, whose value tends to $-\infty$ from the right, contradicting nonnegativity of the integral. Therefore $\sigma_c\le\beta$ (including the possibility $-\infty$), and absolute convergence of the nonnegative transform throughout $\Re s>\beta$ gives the conclusion. The reversed version uses negative residues and a lower bound on $T$.
\end{proof}

\begin{theorem}[Positive-block criterion]\label{thm:block}
Fix a nonempty finite set $S$ of odd primes, positive real weights $b_\ell$, and a fixed $h>0$. For $1/2<\beta<1$, the inequality
\begin{equation}\label{eq:block}
\sum_{\ell\in S} b_\ell\ell\,\overline R_{g_\ell}(v)\le C_\beta e^{(\beta-1)v}\quad(v\text{ sufficiently large})
\end{equation}
implies $\zeta_{K_\ell}(s)\ne0$ for $\Re s>\beta$ and $s\ne1$, for all $\ell\in S$. Under Dirichlet GRH, the same implication follows if \eqref{eq:block} is replaced by the nonresidue deficit bound
\begin{equation}\label{eq:nonres}
-\sum_{\ell\in S} b_\ell\ell\,\overline R_{f_\ell}(v)\le C_\beta e^{(\beta-1)v}.
\end{equation}
If either applicable inequality holds for every $1/2<\beta<1$ (with constants allowed to depend on $S,h,\beta$), all the indicated fields satisfy Dedekind RH. Conversely, their Dedekind RH implies \eqref{eq:block} and \eqref{eq:nonres} for every $\beta>1/2$.
\end{theorem}
\begin{proof}
At any $\rho$ with $\beta<\Re\rho<1$, the sum of ramp transforms in \eqref{eq:block}, before averaging, has residue
\begin{equation}\label{eq:residue}
-\frac{\sum_{\ell\in S} b_\ell\left(\ord_\rho L_\ell+\frac1{\ell-1}\ord_\rho\zeta_{F_\ell}\right)}{\rho(\rho-1)}.
\end{equation}
All orders here are nonnegative integers. There is a pole whenever any factor vanishes. Averaging multiplies the residue by $J_h(\rho)\ne0$. At real $\rho\in(\beta,1)$ the residue is positive. No cancellation is possible, even for coincident nonreal zeros, because the numerator is a positive linear combination of orders. Lemma~\ref{lem:landau} excludes all such poles, and \eqref{eq:fact} proves the conclusion. For \eqref{eq:nonres}, use \eqref{eq:p-f}: Dirichlet GRH removes the $\zeta_{F_\ell}$ terms in this region, and the remaining orders of $L_\ell$ have precisely the same sign as in \eqref{eq:residue}. All factors are already nonzero for $\Re s\ge1$ (with the zeta pole excepted). Let $\beta$ decrease to $1/2$ and use the Dedekind functional equation.

For the converse, effective split-prime Chebotarev under Dedekind RH gives for each fixed field $\pi_{g_\ell}(x)=\delta_g\Li(x)+O_\ell(\sqrt{x}\log x)$, and likewise for the congruence indicator because $\zeta_{F_\ell}$ is a factor without poles in the strip. Thus $\theta_w(x)=\sum_{p\le x}w(p)\log p=\delta_w x+O_\ell(\sqrt{x}\log^2 x)$ for both weights. The identity $W_w(x)=\int_1^x\theta_w(y)y^{-2}\,dy$ follows by finite summation. Integrating the error from $x$ to infinity shows $R_w(\log x)=O_\ell(x^{-1/2}\log^2 x)$: the limiting constant is $c_w$, as follows from \eqref{eq:transform} at $s=1$ or by an elementary Abelian limit of the centered transform. Averaging over a fixed log-window and summing finitely many fixed terms preserves that bound. This is $O_\beta(e^{(\beta-1)v})$.
\end{proof}
\paragraph{Sign discipline.} The source inequality is an \emph{upper} bound for a surplus of non-split quadratic primes, and uses extra twists. Here the unsplit Kummer statistic $R_f$ has a \emph{negative} residue at a real zero of $L_\ell$. Its upper bound alone is not the detector above. We need a lower bound for that surplus, equivalently an upper bound for its deficit, or the upper bound for split primes. This is an essential geometric-transfer obstacle. If one simply substitutes a nonresidue prime weight for the source quadratic weight, the sign is wrong for this proof. The use of positive blocks only works for the correct orientation and with the abelian part removed in the deficit version.

The averaging allows, for instance, $h=\log2$: one bounds the mean of the centered ramp statistic across an octave $x\le y\le2x$. In general a smoothing kernel might introduce zeros that hide a pole. This particular kernel has no such zeros in the open critical strip. We do not claim this averaging principle is unprecedented.

\section{Consequences for Artin's conjecture for 2}\label{sec:artin}
\begin{theorem}[Prime-index hypotheses suffice for density]\label{thm:artin}
Suppose Dedekind RH holds for $K_\ell=\Q(\zeta_\ell,2^{1/\ell})$ for every prime $\ell$. Then
\[
\#\{p\le x:2\text{ is a primitive root modulo }p\}\sim A\,\frac{x}{\log x},\qquad A=\prod_{\ell\ \mathrm{prime}}\left(1-\frac1{\ell(\ell-1)}\right)>0.
\]
In particular, under the exported Dirichlet GRH, \eqref{eq:nonres} for every $\beta>1/2$ on any collection of finite blocks covering the odd primes implies this asymptotic. No squarefree composite-index RH hypothesis is needed for this density conclusion.
\end{theorem}
\begin{proof}
For $p\ne2$, $2$ fails to be a primitive root iff some prime $\ell\mid p-1$ satisfies $2^{(p-1)/\ell}=1$ in $\F_p$. This is exactly complete splitting in $K_\ell$. For a squarefree integer $n$, write $K_n=\Q(\zeta_n,2^{1/n})$. It is the compositum of the $K_\ell$ for $\ell\mid n$: one inclusion is by powers and the other by B\'ezout applied to the integers $n/\ell$, using compatible positive real roots. The cyclotomic field $\Q(\zeta_n)$ is unramified at $2$ (if $n$ is even and squarefree, it equals $\Q(\zeta_{n/2})$). Thus $X^n-2$ is Eisenstein at any prime over $2$, giving $[K_n:\Q]=n\varphi(n)$. For a fixed $Y$, let $N_Y(x)$ count primes $p\le x$ passing every test with $\ell\le Y$. Finite inclusion-exclusion and unconditional Chebotarev for the finitely many $K_n$ give
\[
\lim_{x\to\infty}\frac{N_Y(x)}{\pi(x)}=\prod_{\ell\le Y}\left(1-\frac1{\ell(\ell-1)}\right).
\]
We now bound the tail uniformly, using RH only for the prime-index fields. Write $y=\sqrt{x}/\log^4x$ and $z=\sqrt{x}\log^4x$. The split-prime effective estimate under Dedekind RH (Lang's estimate, or the normalized prime-ideal theorem; see \cite{moree,fan}) is
\[
\pi_{\rm split}(x;K)=\frac{\Li(x)}{[K:\Q]}+O\!\left(\sqrt{x}\left(\log x+\frac{\log|D_K|}{[K:\Q]}\right)\right)
\]
for Galois $K/\Q$, in the present family, with an absolute implied constant. The uniform discriminant bound follows from the tower and the polynomial $X^\ell-2$:
\[
|D_{K_\ell}|\le |D_{F_\ell}|^\ell(\ell^\ell 2^{\ell-1})^{\ell-1},\qquad |D_{F_\ell}|=\ell^{\ell-2},
\]
for $\ell$ odd. Hence $\log |D_{K_\ell}|/[K_\ell:\Q]\ll\log(2\ell)$, and the estimate is $\Li(x)/(\ell(\ell-1))+O(\sqrt x\log x)$ uniformly for $\ell\le x$ (the fixed quadratic case is harmless).

For $Y<\ell\le y$, summing these upper bounds gives
\[
\ll \Li(x)\sum_{\ell>Y}\frac1{\ell(\ell-1)}+y\sqrt{x}\log x
\ll \frac{x}{Y\log x}+\frac{x}{\log^3x}.
\]
For $y<\ell\le z$, disregard the power condition and use Brun--Titchmarsh for $p\equiv1\pmod\ell$. Since $\log(x/\ell)\gg\log x$, the resulting bound is $O((x/\log x)\sum_{y<\ell\le z}1/(\ell-1))=o(x/\log x)$. In detail, the prime reciprocal estimate (or partial summation using $\pi(t)\ll t/\log t$) bounds this sum by $O(1/\log y+\log(\log z/\log y))=O(\log\log x/\log x)$.

For $\ell>z$, if $p$ fails the test then $p\mid 2^m-1$ for $m=(p-1)/\ell\le x/z$. The product of $2^m-1$ over $1\le m\le x/z$ has logarithm $O((x/z)^2)$; the number of distinct prime divisors is at most its logarithm divided by $\log 2$. This gives $O(x/\log^8 x)$ primes. Consequently the difference between $N_Y(x)$ and the desired count is at most $O(x/(Y\log x))+o(x/\log x)$. First take $x\to\infty$, then $Y\to\infty$. The product converges to a positive value because the sum of $1/(\ell(\ell-1))$ converges and all factors are positive. This establishes the asymptotic. The quadratic $K_2$ is covered by Dirichlet GRH, and Theorem~\ref{thm:block} handles every other prime in the stated collection of blocks.
\end{proof}
\paragraph{Limit of the claim.} The argument gives the density asymptotic, not the usual sharper Hooley error term from GRH for all squarefree indices. The finite inclusion-exclusion limit avoids all growing composite indices. No nonabelian RH for any odd $K_\ell$ is established by the Dirichlet theorem alone.

\section{What arithmetic would have to be added}\label{sec:bridge}
We can now state a single precise objective for a proposed extension of the source. Fix $S,h,b_\ell$ as in Theorem~\ref{thm:block} and let
\[
T(v)=-\sum_{\ell\in S}b_\ell\ell\,\overline R_{f_\ell}(v).
\]
Let $U$ be any reference error satisfying the hypotheses of Theorem~\ref{thm:window} (for example, a centered congruence-class ramp error). Suppose that for each sufficiently small $\varepsilon>0$, all large $v$, and every $u\in[t,C_0t]$, $t=A v$, there exist admissible data giving a comparison with $m\ge e^{v/2}$ and positive normalization $D_0$ such that
\begin{equation}\label{eq:newstar}
 m D_0 \bigl(T(v)-T(u)+X U(u)\bigr)\le C_{S,\varepsilon} D_0 e^{k\varepsilon v},\qquad 0\le X\le e^t.
\end{equation}
Here $A\ge1$, $k>0$ are fixed, $C_0$ is sufficiently large as in Corollary~\ref{cor:window}, and crucially $X$ is fixed at the first scale (or otherwise an independently valid common-scale choice is supplied); it cannot vary adversarially with $u$. Data may depend on $u$. This models the original $m\,\mathrm{remainder}\le(\log b)^2\mathrm{budget}$ and the cancellation at two scales.

\begin{proposition}[Exact arithmetic obligation]\label{prop:bridge}
Under the stated uniform hypotheses and Dirichlet GRH, \eqref{eq:newstar} for a fixed block implies Dedekind RH for all its $K_\ell$. If available for a block covering each odd prime, it implies Theorem~\ref{thm:artin}.
\end{proposition}
\begin{proof}
The transforms of $T,U$ satisfy the compact-window hypotheses by Lemma~\ref{lem:transform} and the corresponding centered reference identity. Choose $u\in[t,C_0t]$ with $|T(u)-XU(u)|\le e^{-t}$. Divide \eqref{eq:newstar} by $mD_0>0$ to obtain $T(v)\le C_{S,\varepsilon}e^{(-1/2+k\varepsilon)v}+e^{-Av}$. For every $\beta>1/2$, choose $\varepsilon$ with $k\varepsilon<\beta-1/2$. Theorem~\ref{thm:block} then applies. The sieve is already proved.
\end{proof}
This result is a conditional reduction, not the missing comparison. A program to establish it would require the following independently justified ingredients.
\begin{enumerate}
\item \textbf{Local arithmetic:} create an integral lattice or jet-divisibility statistic whose local contributions discriminate $2$ as an $\ell$th power for primes $p\equiv1\pmod\ell$. Control all other Frobenius types and the ramified primes. Establish the orientation required by \eqref{eq:nonres}. Mere periodic weights cannot achieve this.
\item \textbf{Global comparison:} produce sufficiently many independent short forms and a lower bound on the same lattice's covolume, with matching main heights and residual error $e^{k\varepsilon v}$ after normalization. The square-root saving requires $m\ge e^{v/2}$, not a fixed number of forms. The constructions must accommodate the specified block and all auxiliary $u$ in the compact window.
\item \textbf{Averaging:} if the lattice estimate is pointwise in the original scale, integrate it over an octave only with control of the constants, data changes and normalized errors. It is not legitimate to average unrelated multiplicative covolume inequalities before taking logarithms and normalizing. A valid pointwise bound would suffice without averaging; the weaker averaged statement is a target, not an automatic output.
\item \textbf{Congruences:} if the analogue of the source's fitting argument uses coefficients or twists varying with the second scale, prove uniform selection or modify the inequality. A vector of moments only supplies a common convex combination, not simultaneous pointwise estimates.
\end{enumerate}

\subsection{Exact periodic obstruction for the block}
The source weights are built from congruences and a quadratic character. The Kummer tests cannot be replaced by such periodic functions, even in mean square. Here is an explicit calculation building on the periodic-obstruction perspective of the earlier method-transfer report.
\begin{proposition}\label{prop:periodic}
For the fields for $2$ and a fixed block $S$, define $Z(p)=\sum_{\ell\in S}b_\ell V_\ell(\operatorname{Frob}_p)$ at the good primes. For any real function $a(p)$ of the residue of $p$ modulo a fixed $M$, the prime-density mean square of $Z-a$ is at least $\sum_{\ell\in S} b_\ell^2$, with equality for $a=0$. In particular the best periodic approximation of $\sum b_\ell\ell g_\ell$ has this residual mean square when $M$ is divisible by all $\ell\in S$.
\end{proposition}
\begin{proof}
By the squarefree compositum degree calculation, the Galois group of $K_{\prod S}$ is the direct product of the affine groups. Its commutator subgroup is the product of translations, since for odd $\ell$ the commutator of dilation by $a\ne1$ with a translation gives a nonzero translation. All elements of this commutator subgroup act trivially on the intersection with any abelian cyclotomic extension $\Q(\zeta_M)$. In the compositum with that cyclotomic extension, the Frobenius distribution is uniform on the group by Chebotarev, and conditional on any residue in $(\mathbb Z/M\mathbb Z)^\times$, multiplication by each translation permutes the fiber. The average of $V_\ell$ over a translation coset is zero by its explicit values. Thus $\mathbb E[Z\mid p\bmod M]=0$. The unconditional mean of $V_\ell^2$ is $1$ by the character values; independence of the factors in the direct product gives $\mathbb E[V_\ell V_j]=0$ for $\ell\ne j$. Expand the square to obtain $\mathbb E[(Z-a)^2]=\sum b_\ell^2+\mathbb E[a^2]$. If $M$ contains all the $\ell$, then \eqref{eq:indicator} expresses $\sum b_\ell\ell g_\ell$ as $Z$ plus the periodic function $\sum b_\ell I_\ell$. Subtract it to get the optimal bound. Finitely many bad primes do not affect the mean.
\end{proof}

\section{Further research questions and audit}
\paragraph{Priority 1.} Attack the first nonabelian block $S=\{3\}$, where $K_3$ is the splitting field of $X^3-2$ and $B_3$ is pure cubic. Is there a height or lattice statistic that counts cubic-residue deficits with the sign in \eqref{eq:nonres}? This would imply RH for $\zeta_{K_3}$, not merely a numerical prime race. Before attempting higher rank, verify the local divisibility statement at a split prime and at a nonidentity translation.

\paragraph{Priority 2.} Formalize Theorem~\ref{thm:window} independently of all number-theoretic geometry. Reuse or adapt the source's polar moment argument; the new proof uses gamma-tail estimates and Carath\'eodory. The scalar result needs only continuity, Cauchy estimates, gamma integrals, and the intermediate-value theorem. The vector extension can be a separate theorem.

\paragraph{Priority 3.} Determine whether a uniform family of analytic disks can be obtained for increasing blocks, and whether an arithmetic construction can handle a block aggregate more cheaply than separate fields. A finite-block theorem with constants depending on the block already suffices for the density limit; uniformity is needed for quantitative error estimates or a single infinite-block criterion. Infinite sums cannot be interchanged with meromorphic continuation without additional convergence control.

\paragraph{Priority 4.} Investigate other Kummer bases, and Artin-type generating sets of rank greater than one. Their abelian parts, ramification and entanglement factors need fresh proofs. The case $2$ was chosen because squarefree-index degrees factor exactly, eliminating entanglement in this sieve.

\paragraph{What was verified.} The included Python program tests the character identity for affine groups and for actual primes, verifies the block orthogonality calculations exactly, and illustrates the numerical choice of gamma-window constants. These are finite checks, not a proof of an analytic asymptotic. The arguments are presented in full on paper; no new Lean build or proof is claimed. The archive contains a source map, literature references and a reproducibility log. The strongest open claim, \eqref{eq:newstar}, is deliberately left as a clearly identified obligation. The theorem statements are mathematical reductions, not a resolution of an existing named open conjecture.

\begin{thebibliography}{9}\small
\bibitem{upstream} L. Sernau et al., \emph{GRH for Dirichlet L-functions, in Lean 4}, based on Levent Alp\"oge's proof with Claude, 2026. \url{https://github.com/lsernau/grh-lean}; compact source and guide, tree identifier \nolinkurl{72fead102f663d5c7a407964a13dc0cf9d977207}. See also \url{https://alpo.ge/remark.pdf} (linked by upstream; not independently audited here).
\bibitem{prior} ProveIt, incoming research reports, 25 Nov. 2026. \url{https://github.com/VladimirReshetnikov/ProveIt/tree/main/docs/incoming}. Selected text members inspected; see the companion provenance file.
\bibitem{suzuki} M. Suzuki, On variants of Chebyshev's conjecture, \emph{Ramanujan J.} 68 (2025), article 95, \url{https://doi.org/10.1007/s11139-025-01238-9}. One-sided weighted prime sums and RH/Dirichlet GRH; relevant antecedent.
\bibitem{hooley} C. Hooley, On Artin's conjecture, \emph{J. Reine Angew. Math.} 225 (1967), 209--220. Classical conditional sieve.
\bibitem{moree} P. Moree, Artin's primitive root conjecture---a survey, \emph{Integers} 12 (2012), A13. \url{https://arxiv.org/abs/math/0412262}. Section on Hooley's method and the uniform split-prime estimate.
\bibitem{fan} S. Fan and P. Pollack, Counting primes with a given primitive root, uniformly, \emph{Mathematika} 71 (2025), e70055, \url{https://doi.org/10.1112/mtk.70055}; \url{https://arxiv.org/abs/2505.05601}.
\bibitem{mv} H. L. Montgomery and R. C. Vaughan, \emph{Multiplicative Number Theory I: Classical Theory}, Cambridge, 2007. Landau's theorem and analytic prime-number methods.
\bibitem{neukirch} J. Neukirch, \emph{Algebraic Number Theory}, Springer, 1999. Artin formalism, abelian class field theory, discriminants, Chebotarev and zeta functions.
\end{thebibliography}
\end{document}
